\chapter{Theoretische Grundlagen}
\label{chapter:kapitel2}

Erarbeitung der wissenschaftlichen Basis zu Usability, den Interaktionsprinzipien der ISO 9241-110 und der heuristischen Evaluation als 
Untersuchungsmethode.

\section{Usability als Qualitätsmerkmal interaktiver Systeme}
\label{section:usability}

Der Begriff Usability, der im Deutschen auch als Gebrauchstauglichkeit oder Benutzbarkeit bezeichnet wird, beschreibt die 
Eignung eines Systems, von seinen vorgesehenen Nutzern zur Erreichung ihrer Ziele eingesetzt werden zu 
können.\footcite[Vgl.][S. 3]{richter_2010}

Dabei greift das Verständnis, das Usability allein auf die Gestaltung der Benutzeroberfläche reduziert, zu kurz. Qualitätskriterien 
wie die Anordnung von Bedienelementen oder die Anzahl notwendiger Klicks betreffen lediglich die 
Oberflächenebene.\footcite[Vgl.][S. 3]{richter_2010} In einem weiteren Verständnis muss die Benutzbarkeit eines Systems 
stets im Kontext seiner konkreten Verwendung beurteilt werden. Ein System weist demnach eine hohe Usability auf, wenn es von den 
vorgesehenen Nutzern einfach erlernt und effizient verwendet werden kann und diese ihre beabsichtigten Ziele zufriedenstellend 
erreichen.\footcite[Vgl.][S. 3]{richter_2010} Dass die Gebrauchstauglichkeit nicht ausschließlich eine Eigenschaft des 
Produktes selbst ist, verdeutlicht das Verhältnis von Werkzeug und Aufgabe. Dasselbe Werkzeug kann für eine Aufgabe gut, für eine 
andere hingegen schlecht geeignet sein. Usability bezeichnet somit, wie gut Nutzer ein Werkzeug in ihrem Umfeld zur Bewältigung ihrer 
Aufgaben einsetzen können.\footcite[Vgl.][S. 5]{richter_2010}

Der Usability-Begriff ist eng mit dem Forschungsfeld der \ac{HCI} (im Deutschen als Mensch-Computer-Interaktion bezeichnet) verbunden, 
das sich mit der Gestaltung von Systemen befasst, die Menschen bei der produktiven und sicheren Ausführung ihrer Tätigkeiten 
unterstützen.\footcite[Vgl.][S. 23]{issa_sustainable_2022} Während HCI die gestalterischen und interaktiven Aspekte eines Systems 
adressiert, dient Usability als Maßstab dafür, ob ein System effizient, effektiv, sicher und erlernbar ist und den Nutzern 
Zufriedenheit bei der Aufgabenerfüllung verschafft.\footcite[Vgl.][S. 4]{issa_sustainable_2022} Eine zentrale Funktion erfüllt die 
Usability dabei im Entwicklungsprozess. Durch eine systematische Usability-Bewertung lassen sich Stärken und Schwächen eines Systems 
bereits vor dessen Auslieferung erkennen und notwendige Anpassungen vornehmen.\footcite[Vgl.][S. 4]{issa_sustainable_2022} Damit zeigt 
sich, dass Gebrauchstauglichkeit nicht als eindimensionale Eigenschaft, sondern als ein aus mehreren Qualitätsmerkmalen zusammengesetztes 
Konstrukt zu verstehen ist.

Die praktische Bedeutung der Gebrauchstauglichkeit zeigt sich besonders an den Folgen ihres Fehlens. Eine unzureichende Usability äußert
sich unter anderem darin, dass Mitarbeiter langsamer als erwartet mit einem System arbeiten, die Einarbeitung viel Zeit beansprucht, 
Prozessvorgaben umgangen werden oder wiederkehrende Bedienfehler Schäden verursachen.\footcite[Vgl.][S. 7]{richter_2010}


\section{Interaktionsprinzipien nach ISO 9241-110}
\label{grundsaetzeISO}

Während Abschnitt~\ref{section:usability} die Gebrauchstauglichkeit als übergeordnetes Qualitätsmerkmal beschrieben hat, konkretisiert die 
Norm ISO 9241-110:2020 diese durch sieben Interaktionsprinzipien, die für die Gestaltung und Bewertung interaktiver Systeme als wesentlich 
bestimmt wurden.\footcite[Vgl.][Abschnitt 4.1]{DIN_ISO_9241-110:2020} Die Norm beschreibt diese Prinzipien als technologieunabhängige 
und generische Kriterien, die unabhängig von einer bestimmten Interaktionstechnik bei der Analyse, Gestaltung und Bewertung interaktiver 
Systeme angewendet werden können.\footcite[Vgl.][Einleitung]{DIN_ISO_9241-110:2020} Ihre Reihenfolge im Dokument stellt dabei ausdrücklich 
keine Rangfolge dar. Vielmehr hängt die Bedeutung des einzelnen Prinzips vom jeweiligen Nutzungskontext, den Benutzergruppen und der gewählten 
Interaktionstechnik ab.\footcite[Vgl.][Abschnitt 4.5]{DIN_ISO_9241-110:2020}

Die Interaktionsprinzipien lösen die in der Vorgängerfassung von 2006 als „Grundsätze der Dialoggestaltung" bezeichneten Kriterien ab. Im Zuge 
der Überarbeitung wurde das Prinzip der Individualisierung in die Steuerbarkeit integriert und mit der Benutzerbindung ein neues, siebtes Prinzip 
eingeführt.\footcite[Vgl.][Nationales Vorwort, Änderungen]{DIN_ISO_9241-110:2020} Die terminologische Verschiebung von der „Dialoggestaltung" hin zu 
„Interaktionsprinzipien" verdeutlicht zugleich den erweiterten Anwendungsbereich der Norm, der nicht mehr nur den klassischen Bildschirmdialog, 
sondern die Interaktion mit interaktiven Systemen im weiteren Sinne umfasst.

Die Norm ordnet jedem Prinzip Hauptkategorien von allgemeinen Gestaltungsempfehlungen zu, anhand derer sich das jeweilige Prinzip konkretisieren 
lässt. Diese Zuordnung ist in Tabelle~\ref{tab:interaktionsprinzipien} zusammengefasst.

\newpage
\begin{table}[htbp]
	\centering
	\ra{1.9}
	\caption[Interaktionsprinzipien und Hauptkategorien der entsprechenden Empfehlungen]{Interaktionsprinzipien und Hauptkategorien der entsprechenden Empfehlungen\footcite[Enthalten in:][Tabelle 1]{DIN_ISO_9241-110:2020}}
	\label{tab:interaktionsprinzipien}
	\begin{tabularx}{0.98\textwidth}{|l|>{\RaggedRight}X|}
		\hline
		\textbf{Interaktionsprinzip} & \textbf{Kategorien von Empfehlungen} \\
		\hline
		Aufgabenangemessenheit &
		a) Identifizierung der Angemessenheit des interaktiven Systems für eine bestimmte Aufgabe \newline
		b) Aufwandsoptimierung bei der Aufgabenerledigung \newline
		c) Standardauswahlmöglichkeiten zur Aufgabenunterstützung \\
		\hline
		Selbstbeschreibungs\-fähigkeit &
		a) Vorhandensein und Offensichtlichkeit der Informationen \newline
		b) eindeutige Anzeige des Systemzustands \\
		\hline
		Erwartungskonformität &
		a) angemessene(s) Systemverhalten und -reaktionen \newline
		b) Konsistenz (interne und externe) \newline
		c) Änderungen im Nutzungskontext \\
		\hline
		Erlernbarkeit &
		a) Entdeckung (von Informationen und Steuerelementen, nach denen Benutzer suchen) \newline
		b) Exploration (von Informationen und Steuerelementen, die Benutzer aufgefunden haben) \newline
		c) Retention (von Informationen über das System) \\
		\hline
		Steuerbarkeit &
		a) Unterbrechung durch den Benutzer \newline
		b) Flexibilität \newline
		c) Individualisierung \\
		\hline
		Robustheit gegen Benutzungsfehler &
		a) Vermeidung von Benutzungsfehlern \newline
		b) Toleranz gegen Benutzungsfehler \newline
		c) Behebung von Benutzungsfehlern \\
		\hline
		Benutzerbindung &
		a) Motivieren des Benutzers \newline
		b) Vertrauenswürdigkeit des Systems \newline
		c) Zunehmende Einbeziehung des Benutzers in das System \\
		\hline
	\end{tabularx}
\end{table}

Im Folgenden werden die sieben Interaktionsprinzipien einzeln erläutert.

Die \textbf{Aufgabenangemessenheit} beschreibt, inwieweit ein interaktives System die Benutzer bei der Erledigung ihrer Aufgaben unterstützt. 
Sie ist gegeben, wenn die Bedienfunktionen und die Benutzer-System-Interaktionen auf den charakteristischen Eigenschaften der Aufgabe beruhen 
und nicht auf der zur Erfüllung der Aufgabe eingesetzten Technologie.\footcite[Vgl.][Abschnitt 5.1.1]{DIN_ISO_9241-110:2020} Die Norm setzt 
dabei voraus, dass die Aufgaben selbst auf den Erfordernissen der Benutzer basieren, und betont, dass dem Benutzer keine Schritte auferlegt 
werden sollten, die nicht aus der Aufgabe, sondern allein aus der Technologie resultieren.\footcite[Vgl.][Abschnitt 5.1.1]{DIN_ISO_9241-110:2020}

Die \textbf{Selbstbeschreibungsfähigkeit} ist gegeben, wenn das System dem Benutzer bei Bedarf angemessene Informationen bereitstellt, die seine 
Fähigkeiten und seine Nutzung unmittelbar offensichtlich machen, ohne dass hierzu unnötige Interaktionen erforderlich 
werden.\footcite[Vgl.][Abschnitt 5.2.1]{DIN_ISO_9241-110:2020} Ein System ist in dem Maße selbstbeschreibungsfähig, in dem für den Benutzer 
jederzeit erkennbar ist, an welcher Stelle des Systems er sich befindet, welche Handlungen möglich sind und wie diese ausgeführt werden 
können.\footcite[Vgl.][Abschnitt 5.2.1]{DIN_ISO_9241-110:2020}

Die \textbf{Erwartungskonformität} liegt vor, wenn sich das Verhalten des Systems auf Grundlage des Nutzungskontexts und allgemein anerkannter 
Konventionen vorhersehen lässt.\footcite[Vgl.][Abschnitt 5.3.1]{DIN_ISO_9241-110:2020} Die Norm unterscheidet hierbei zwischen interner 
Konsistenz, also der Einheitlichkeit innerhalb des Systems selbst, und externer Konsistenz, die sich auf die Übereinstimmung mit anderen, 
vergleichbaren Systemen bezieht, mit denen der Benutzer voraussichtlich interagiert.\footcite[Vgl.][Abschnitt 5.3.1]{DIN_ISO_9241-110:2020}

Die \textbf{Erlernbarkeit} umfasst die Unterstützung der Benutzer beim Entdecken und Explorieren der Systemfunktionen sowie die Minimierung des 
Lernaufwands und die Bereitstellung von Unterstützung, wenn Lernen erforderlich ist.\footcite[Vgl.][Abschnitt 5.4.1]{DIN_ISO_9241-110:2020} 
Die Norm gliedert dieses Prinzip in die Aspekte der Entdeckung, der Exploration und der Retention, also des Wiederauffindens bereits erlernter 
Informationen über das System.\footcite[Vgl.][Abschnitt 5.4.1]{DIN_ISO_9241-110:2020}

Die \textbf{Steuerbarkeit} erlaubt es dem Benutzer, die Kontrolle über die Benutzungsschnittstelle und die Interaktionen zu behalten, 
einschließlich Geschwindigkeit, Abfolge und Individualisierung der 
Benutzer-System-Interaktion.\footcite[Vgl.][Abschnitt 5.5.1]{DIN_ISO_9241-110:2020} Während sich die Flexibilität auf alternative Möglichkeiten 
bezieht, die das System dem Benutzer unmittelbar bereitstellt, betrifft die Individualisierung die Fähigkeit des Systems, sich an die 
individuellen Kenntnisse, Fertigkeiten und Vorlieben des Benutzers anpassen zu lassen.\footcite[Vgl.][Abschnitt 5.5.1]{DIN_ISO_9241-110:2020}

Die \textbf{Robustheit gegen Benutzungsfehler} bezeichnet die Unterstützung des Benutzers beim Vermeiden von Fehlern, die Toleranz gegenüber 
erkennbaren Fehlern sowie die Unterstützung bei der Fehlerbehebung.\footcite[Vgl.][Abschnitt 5.6.1]{DIN_ISO_9241-110:2020} Fehlertoleranz liegt 
nach der Norm dann vor, wenn das beabsichtigte Ergebnis trotz erkennbar fehlerhafter Eingaben entweder ohne oder mit nur minimalem 
Korrekturaufwand seitens des Benutzers erreicht werden kann.\footcite[Vgl.][Abschnitt 5.6.1]{DIN_ISO_9241-110:2020}

Die \textbf{Benutzerbindung} schließlich beschreibt, inwieweit das System Funktionen und Informationen auf einladende und motivierende Weise 
darstellt und so eine kontinuierliche Interaktion fördert.\footcite[Vgl.][Abschnitt 5.7.1]{DIN_ISO_9241-110:2020} Als neu eingeführtes Prinzip 
betont die Norm, dass Benutzerbindung auf ethisch korrekte Weise zu erfolgen hat und das Einhalten der übrigen sechs Prinzipien eine wesentliche 
Voraussetzung für ihre Wirksamkeit darstellt.\footcite[Vgl.][Abschnitt 5.7.1]{DIN_ISO_9241-110:2020}

Die Norm weist darauf hin, dass die Prinzipien nicht streng unabhängig voneinander sind und sich inhaltlich überschneiden können, sodass sich ein 
einzelnes Usability-Problem auf mehrere Prinzipien zugleich beziehen kann.\footcite[Vgl.][Abschnitt 4.5]{DIN_ISO_9241-110:2020} Da die Prinzipien 
selbst nicht unmittelbar überprüfbar sind, stellt die Norm für jedes Prinzip konkrete Gestaltungsempfehlungen bereit, deren Anwendung jedoch 
nicht bedeutet, dass das übergeordnete Prinzip in Gänze erfüllt ist.\footcite[Vgl.][Abschnitt 4.1]{DIN_ISO_9241-110:2020} Diese 
Gestaltungsempfehlungen bilden die Grundlage für die in Kapitel~\ref{section:untersuchungsdesign} entwickelte Operationalisierung zu einem 
Kriterienkatalog, der als Bewertungsmaßstab der heuristischen Evaluation dient.


\section{Heuristische Evaluation als Methode}
\label{section:heuristische}

Die heuristische Evaluation ist eine Methode der expertenbasierten Usability-Bewertung,
bei der eine kleine Gruppe von Gutachtern eine Benutzungsschnittstelle eigenständig prüft
und deren Übereinstimmung mit einer Menge anerkannter Usability-Prinzipien, den
sogenannten Heuristiken, bewertet.\footcite[Vgl.][S.~25]{nielsen1994heuristic} Begründet wurde
sie von Molich und Nielsen, die die Komplexität umfangreicher Gestaltungsrichtlinien
bewusst auf eine überschaubare Liste weniger Prinzipien reduzierten, um die Methode für die
praktische Anwendung handhabbar zu machen.\footcite[Vgl.][S. 338]{molich1990improving} Als
Inspektionsmethode kommt die heuristische Evaluation ohne die Beteiligung tatsächlicher
Nutzer aus. Anders als beim nutzerbasierten Usability-Test beurteilen Gutachter
stellvertretend, ob eine Oberfläche etablierten Gestaltungsprinzipien
genügt.\footcite[Vgl.][S. 249]{nielsen1990heuristic} Sie gilt damit als
ressourcenökonomisches Verfahren, das bereits in frühen Phasen des Entwicklungsprozesses
eingesetzt werden kann.\footcite[Vgl.][S. 255]{nielsen1990heuristic}

Im Ablauf durchläuft jeder Gutachter die Benutzungsschnittstelle mehrfach und gleicht ihre
einzelnen Dialogelemente systematisch mit den Heuristiken ab. Damit die Bewertungen
unvoreingenommen bleiben, prüfen die Gutachter zunächst unabhängig voneinander und erst nach
Abschluss aller Einzelprüfungen werden die Befunde
zusammengeführt.\footcite[Vgl.][S.~25]{nielsen1994heuristic} Das Ergebnis ist eine Liste der
identifizierten Usability-Probleme, in der jedes Problem auf dasjenige Prinzip zurückgeführt
wird, gegen das die Gestaltung verstößt.\footcite[Vgl.][S.~25]{nielsen1994heuristic} Die gefundenen
Probleme werden anschließend nach ihrem Schweregrad eingeordnet. Bereits in den
Gründungsarbeiten wurden solche Probleme, die Nutzer an der sinnvollen Erledigung ihrer
Aufgaben hindern können, gesondert als schwerwiegend
ausgewiesen.\footcite[Vgl.][S. 340]{molich1990improving}

Ein zentraler empirischer Befund der Methode betrifft die Anzahl der Gutachter. Einzelne
Gutachter decken je nach untersuchtem System nur etwa 20 bis 51\,\% der vorhandenen
Usability-Probleme auf.\footcite[Vgl.][S. 249]{nielsen1990heuristic} Auch unter erfahrenen
Fachleuten zeigt sich dieses Problem. In einer Befragung von 77 Entwicklern wurden im
Schnitt lediglich 37\,\% der in einem Dialog enthaltenen Probleme
erkannt.\footcite[Vgl.][S. 340]{molich1990improving} Da unterschiedliche Gutachter jedoch
unterschiedliche Probleme aufdecken, lässt sich die Trefferquote durch das Zusammenführen
mehrerer unabhängiger Bewertungen deutlich
erhöhen.\footcite[Vgl.][S. 253]{nielsen1990heuristic} Nielsen und Molich empfehlen daher den
Einsatz von drei bis fünf Gutachtern.\footcite[Vgl.][S. 254]{nielsen1990heuristic} Auch
aktuellere Untersuchungen bestätigen, dass Gutachter bei der Bewertung desselben Systems
trotz identischer Vorgaben zu unterschiedlichen Ergebnissen gelangen, und empfehlen einen
Ablauf, der individuelle Einzelbewertungen mit einer anschließenden gemeinsamen Validierung
verbindet.\footcite[Vgl.][S. 363]{inal2021collaborative}

Die Stärken der Methode liegen in ihrer Wirtschaftlichkeit und einfachen Anwendbarkeit. Sie
ist kostengünstig, intuitiv, erfordert keine aufwendige Vorbereitung und ist früh im
Entwicklungsprozess einsetzbar.\footcite[Vgl.][S. 254]{nielsen1990heuristic} Zugleich
bestehen methodische Grenzen. So hängen die Ergebnisse von  der Anzahl und Expertise der Gutachter
ab. Ein weiteres Problem besteht darin, dass sie auf Fehler hinweist, ohne unmittelbare Lösungsvorschläge zu liefern,
zudem sind die Ergebnisse durch die Perspektive der jeweiligen Gutachtenden 
geprägt.\footcite[Vgl.][S. 254]{nielsen1990heuristic} Diese Abhängigkeit von der
Gutachterkonstellation ist für die vorliegende Arbeit von besonderer Bedeutung, da die
untersuchte Anwendung von einer einzelnen Person betrachtet wird. Die daraus resultierenden
Einschränkungen werden in Kapitel \ref{grenzen} vertieft.

Klassischerweise stützt sich die heuristische Evaluation auf die von Molich und Nielsen
formulierten Heuristiken.\footcite[Vgl.][S. 338]{molich1990improving} Für die vorliegende
Untersuchung treten an deren Stelle die Interaktionsprinzipien der ISO 9241-110, die als
Heuristiken im Sinne der Methode herangezogen werden. Dieses Vorgehen ist insofern
naheliegend, als die Norm ihre allgemeinen Gestaltungsempfehlungen ausdrücklich unter
anderem aus den Arbeiten von Molich und Nielsen
ableitet\footcite[Vgl.][Einleitung]{DIN_ISO_9241-110:2020} und die Interaktionsprinzipien
unmittelbar auf die Verbesserung der Gebrauchstauglichkeit
zielen.\footcite[Vgl.][Abschnitt 4.4]{DIN_ISO_9241-110:2020} Die konkrete Überführung der
Interaktionsprinzipien in einen Kriterienkatalog als Bewertungsmaßstab erfolgt in Kapitel
\ref{section:untersuchungsdesign}.
